\section{Why analytic local factoriality is different}
\label{sec:local-convention}

Global Weil-divisor $\Q$-factoriality does not require every analytic local
ring to be $\Q$-factorial.  The distinction already occurs for projective
threefolds with one ordinary double point; see
\cite[Definition~5.1, Remark~(b)]{Reid}.  The following product example
explains why the stronger local requirement cannot be inserted into
Theorem~\ref{thm:finite}.

\begin{proposition}\label{prop:local-convention}
There is a smooth projective fourfold $X$ with pseudo-effective, non-nef
canonical divisor for which no ordinary $K_X$-negative program can reach a
nef model while keeping every model analytically locally
$\Q$-factorial.  Its unique possible first negative birational contraction
is divisorial.  Its target is nevertheless globally factorial and admits
a nef canonical divisor.
\end{proposition}

\begin{proof}
\emph{Construction.}
Let $b:B\to\Pbb^4$ be the blowup of a point, let
$H=b^*\cO_{\Pbb^4}(1)$, and let $F$ be the exceptional divisor.
The graph of projection from that point embeds $B$ in
$\Pbb^4\times\Pbb^3$, and the restriction of $\cO(1,1)$ is $2H-F$.
Thus $2H-F$ is very ample.  Since $H$ is globally generated,
\[
 6H-2F=2(2H-F)+2H
\]
is very ample as well.  Take a general smooth member
$U\in|6H-2F|$.
Its image $U_0\subset\Pbb^4$ is smooth away from the blown-up point.
The quadratic jets of the defining sextics at that point are arbitrary,
so a general member has a nondegenerate quadratic initial term.
The holomorphic Morse lemma identifies its unique singularity with
\[
 xy-zw=0.
\]
The restriction $\pi:U\to U_0$ resolves this node, with exceptional
quadric
\[
 E=F|_U\simeq\Pbb^1\times\Pbb^1,
 \qquad \cO_U(E)|_E=\cO_E(-1,-1).
\]
Adjunction gives
\begin{equation}\label{eq:local-example-canonical}
 K_U=H|_U+E,
\end{equation}
since $K_B=-5H+3F$.

Integral weak Lefschetz for the smooth ample threefold $U\subset B$
\cite[Section~7, Corollary~7.3]{MilnorMorse1963} gives
\[
 H^1(U,\Z)=0,
 \qquad H^2(U,\Z)=\Z[H|_U]\oplus\Z[E].
\]
This is the blowup calculation underlying the classical example in
\cite[Definition~5.1, Remark~(b)]{Reid}.
Let $C$ be a smooth elliptic curve and set $X=U\times C$.
Its canonical divisor is the pullback of \eqref{eq:local-example-canonical},
so is effective.  If $\ell$ is either ruling in $E\times\{c\}$, then
$K_X\cdot\ell=-1$; in particular $K_X$ is not nef.

\emph{All negative contractions have the same target.}
Put $D=E\times C$ and let $h=\{e\}\times C$ be a horizontal curve.
A curve not contained in $D$ has nonnegative degree against both the
effective divisor $D$ and the nef hyperplane pullback.  Thus every
$K_X$-negative curve lies in $D$.
The two rulings of $E$ have the same class in $H_2(U,\R)$: both have degrees
$0,-1$ against the displayed basis of $H^2(U,\R)$.
Künneth and $H^1(U,\R)=0$ now show that every integral curve
$\gamma\subset D$ has class
\begin{equation}\label{eq:local-example-curve}
 [\gamma]=a[\ell]+b[h],\qquad a,b\geq0.
\end{equation}
Here $a$ is the sum of the projection degrees to the two factors of
$E=\Pbb^1\times\Pbb^1$, and $b$ is the projection degree to $C$.
Moreover $K_X\cdot\gamma=-a$, so negativity means $a>0$.

Let $f:X\to Z$ be any nontrivial projective birational contraction with
connected fibres, normal compact Kähler target, and $-K_X$ relatively
ample.  Choose a Kähler class $\omega_Z$ and set $\alpha=f^*\omega_Z$.
For a curve, vanishing of its degree against $\alpha$ is equivalent to
being contracted by $f$.
In particular $\alpha\cdot\ell\geq0$ and $\alpha\cdot h>0$: the
horizontal curve cannot be contracted because $K_X\cdot h=0$.
A contracted curve $\gamma$ has the form
\eqref{eq:local-example-curve} with $a>0$, and
\[
 0=\alpha\cdot\gamma
   =a(\alpha\cdot\ell)+b(\alpha\cdot h).
\]
It follows that $b=0$ and $\alpha\cdot\ell=0$.
Thus all rulings in every $E\times\{c\}$ are contracted, and $f$ is not
small.

Set $q=\pi\times\id_C:X\to U_0\times C$.
Its nontrivial fibres are the quadrics $E\times\{c\}$, which are connected
by their rulings, so $f$ is constant on every $q$-fibre.
Conversely, every curve in an $f$-fibre has the form
\eqref{eq:local-example-curve} with $b=0$, and is contracted by $q$.
A connected projective fibre is connected by chains of curves, so $q$ is
constant on every $f$-fibre.  The two targets are isomorphic: the reduced
image of $(q,f)$ in their product projects properly, bijectively and
bimeromorphically to each normal target, hence finitely and isomorphically.
Therefore $q$ is the only possible first negative birational contraction.
It is divisorial, and its relative numerical rank is one because both
quadric rulings have the same nonzero numerical class.

\emph{Failure of analytic local $\Q$-factoriality.}
Near the singular locus of $U_0\times C$, consider the prime analytic
divisor
\[
 P=\{x=z=0\}\times\text{(a disk)}
   \ \subset\ \{xy-zw=0\}\times\text{(a disk)}.
\]
On the blowup of the singular axis, its strict transform meets each
exceptional quadric in a ruling line.  If $nP$ were Cartier near an axis
point for some integer $n>0$, its pullback would have the form
$n\widetilde P+kE_{\mathrm{exc}}$.
The associated line bundle restricts trivially to the quadric over that
point, because it is pulled back from a line bundle at a point.
On the other hand its restriction is
\[
 \cO_{\Pbb^1\times\Pbb^1}(n-k,-k),
\]
up to interchanging the factors.  Triviality forces $k=0$ and $n=0$, a
contradiction.  The target is not analytically locally $\Q$-factorial.
Since $K_X$ is not nef and there is no small negative contraction,
there is no permissible first step under that local convention.

\emph{Compatibility with the global convention.}
The same example has a valid one-step global program.  Indeed,
$-E$ is $\pi$-ample, so $-K_X$ is $q$-ample, and adjunction on $U_0$ gives
\[
 \cO_{U_0\times C}(K_{U_0\times C})
 \simeq\operatorname{pr}_1^*\cO_{U_0}(1),
\]
which is nef.  On the blowup of the node the exceptional divisor has log
discrepancy two; its smooth product with $C$ is a log resolution with the
same log discrepancy.  Hence the target is terminal, in particular klt.

It remains to check that the local divisor $P$ is not a global
factoriality obstruction.  The displayed integral cohomology of $U$ and
Hodge decomposition give $H^1(U,\cO_U)=H^2(U,\cO_U)=0$ and
$\Pic(U)=\Z[H|_U]\oplus\Z[E]$.
The exponential sequence and Künneth therefore give
\[
 \Pic(X)=\Z[H]\oplus\Z[D]\oplus\operatorname{pr}_C^*\Pic(C).
\]
Strict transform and divisor pushforward identify
\[
 \Cl(U_0\times C)=\Pic(X)/\Z[D].
\]
Every surviving class descends to a Cartier class, either from the
hyperplane bundle on $U_0$ or from a line bundle on $C$.
Chow's theorem makes every global analytic Weil divisor on this projective
target algebraic, so every such divisor is Cartier
\cite[Section~4, no.~19, Proposition~13]{SerreGAGA1956}.
GAGA identifies its global coherent
rank-one reflexive sheaves with algebraic divisorial sheaves, which are
therefore invertible as well
\cite[Section~3, no.~12, Theorems~2--3]{SerreGAGA1956}. The plane $P$ is defined only on an analytic
neighbourhood; it need not extend to a divisor on the whole compact space.
\end{proof}
